\section{Experiments}\label{sec:experiments}

\subsection{Evaluation setups}\label{sec:eval-setup}

\begin{figure}[!t]
\centering\footnotesize
\adjustbox{max width=\columnwidth}{%
\begin{tabular}{@{}p{1.15\columnwidth}@{}}
\toprule
\multicolumn{1}{c}{\textbf{PaLM outputs on NQ (factoid statements)}} \\
Millie Inbetween is a British comedy television series. It premiered on 24 February 2014 on BBC One. The first series was produced by John Yorke and Phil Clymer. \\ \midrule
\multicolumn{1}{c}{\textbf{PaLM outputs on SQA (reasoning chains)}} \\
The highest point of Mount Wycheproof is 70 metres. Edmund Hillary climbed Mount Everest, which is 8,848 metres. So Mount Wycheproof would be a breeze for Edmund Hillary. \\ \midrule
\multicolumn{1}{c}{\textbf{LaMDA outputs on QReCC (knowledge-intensive dialogs)}} \\
\begin{tabular}{@{}ll@{}}
When was Welsh social reformer Robert Owen born?
& \multirow{4}{*}{
$\left.\vphantom{\rule{0pt}{5ex}}\right\}$ context
} \\
Robert Owen was born on 14 May 1771 \\
\dots \\
Did he have another job? \\
\end{tabular} \\
In 1810 he moved to Manchester and established a draper's shop. \\
\bottomrule
\end{tabular}
}
\caption{
\label{fig:datasets}
\textbf{Examples of input passages.} For QReCC, prior dialog turns are also given as the context.
}
\end{figure}

\approach{} aspires to be a general-purpose method for improving the attribution of any text generation model in any text domain.
We thus construct evaluation benchmarks by taking the task input from three diverse datasets,
and prompting different generation models to produce \emph{long-form outputs} which may contain ``hallucinations,''
as demonstrated in Figure~\ref{fig:datasets}.
These long-form outputs serve as input text passages to \approach{}.
We generate 150 development and 150 test passages for each combination of generation model and source dataset. %






\paragraph{Factoid statements}
We prompt PaLM 540B and GPT-3 text-davinci-002 to generate long-form answers to questions from the Natural Questions dev set \cite[NQ;][]{Kwiatkowski2019NaturalQA}. The resulting passages are mostly coherent but often contain factual errors. This setup examines the ability to attribute a diverse range of factoid knowledge.

\paragraph{Reasoning chains}
Language models can generate reasoning chains to answer complex questions \cite{Wei2022COT}. We use PaLM and GPT-3 to generate reasoning chains for the StrategyQA train set \cite[SQA;][]{Geva2021DidAU}.
This setup tests whether the revision model can provide better attribution for intermediate steps of reasoning, while preserving the overall reasoning process.

\paragraph{Knowledge-intensive dialogs}
We consider the conversational QA task from the QReCC dev set \cite{Anantha2021QReCCQ}.
Given the previous dialog turns, which are rounds of questions and answers ($Q_1, A_1, Q_2, A_2, \dots, Q_k$), we use LaMDA and GPT-3 to answer %
to the final question $Q_k$ conditioned on the dialog history.
The answer tends to be context-dependent, featuring pronouns and implicit references. All dialog turns are given alongside the answer as inputs to the revision model.




\subsection{Models}\label{sec:baselines}
We compare \approach{} to several systems that have a research-and-revise workflow.

\paragraph{EFEC}
We consider EFEC~\cite{thorne2021evidence} as a representative fine-tuned editor.
EFEC fine-tunes a T5-based model to revise text conditioned on multiple evidence snippets using both semi-supervised and fully-supervised approaches.
We compare against their fully-supervised approach, which performed best in their experiments.
EFEC uses a neural retrieval model~\cite{karpukhin2020dpr} to retrieve from Wikipedia; however, not all passages in our experiments are supported by Wikipedia articles.
To more fairly compare the editing capabilities of EFEC, we instead use the evidence retrieved by our research stages (CQGen and web search). %
Note that the EFEC editor conditions on multiple pieces of evidence at once, while our editor iteratively conditions on one at a time.


\paragraph{LaMDA}
LaMDA~\cite{Thoppilan2022LaMDA} generates responses in three steps: 
1) generate a ``base response'';
2) generate search queries from the base response;
3) generate a ``revised response'' conditioned on the base response and retrieved evidence.
To apply LaMDA on a given text $x$, we simply set the base response in step 1 to $x$, and then run steps 2 and 3 (we call these latter two stages ``LaMDA Research'').
LaMDA was trained as a dialog system, and always expects a dialog context where the user speaks first.
So, for non-dialog tasks, we insert an artificial user utterance as dialog history: \emph{``Tell me something interesting.''}
For the attribution report, we take all evidence documents retrieved by LaMDA during its research process.

\paragraph{\approach{}}
Our model uses few-shot prompting on PaLM 540B for query generation, the agreement model, and the edit model.
We use the same prompts for all tasks except when the context comes from a dialog,
    where we slightly modify the prompts to use the dialog context (e.g., CQGen now maps dialog context + $x$ to queries).
The query-evidence relevance model $S_\text{relevance}$ is a pretrained T5-large model \cite{raffel20t5} fine-tuned following \citet{ni2021gtr} on MS MARCO \cite{nguyen2016msmarco}.
See Appendix~\ref{app:model-details} for the few-shot prompting strategies and more modeling details.

\newcommand{\mainTableDatasetName}[1]{\multicolumn{9}{c}{\textbf{#1}}}

\begin{table}[!t]
\setlength{\tabcolsep}{4pt}
\centering
\adjustbox{max width=\columnwidth}{%
\begin{tabular}{
@{}l
r@{ $\to$ }r
r@{ $\to$ }r
rrrr
@{}}
\toprule
 & \multicolumn{4}{c}{Attribution} & \multicolumn{3}{c}{Preservation}
 \\
 \cmidrule(lr){2-5} \cmidrule(lr){6-8}
 Model
 & \multicolumn{2}{c}{auto-AIS} & \multicolumn{2}{c}{AIS} & \MN{intent} & \MN{Lev} & \MN{comb} & $\HM{}$
 \\
\midrule
\mainTableDatasetName{PaLM outputs on NQ} \\
EFEC
& 45.6 & 64.3 & 35.4 & 48.3 & 16.0 & 39.1 & 10.4 & 17.1 \\
LaMDA
& 39.5 & 49.9 & 18.3 & 30.4 & 26.0 & 39.6 & 21.1 & 24.9 \\
\approach{}
& 45.6 & 54.9 & 35.4 & 43.4 & 90.0 & 89.6 & 83.1 & \textbf{57.0} \\ \midrule
\mainTableDatasetName{PaLM outputs on SQA} \\
EFEC
& 37.8 & 58.6 & 24.5 & 51.7 & 6.0 & 31.0 & 3.8 & 7.1 \\
LaMDA
& 32.7 & 43.2 & 15.8 & 27.0 & 40.0 & 46.4 & 33.7 & 30.0 \\
\approach{}
& 37.6 & 45.1 & 24.5 & 31.5 & 92.6 & 89.9 & 84.6 & \textbf{45.9} \\ \midrule
\mainTableDatasetName{LaMDA outputs on QReCC} \\
EFEC
& 19.1 & 47.4 & 13.2 & 48.7 & 39.7 & 39.4 & 23.7 & 31.9 \\
LaMDA
& 16.4 & 36.2 & 16.0 & 27.1 & 21.3 & 24.8 & 12.0 & 16.6 \\
\approach{}
& 18.8 & 29.4 & 13.2 & 28.3 & 95.6 & 80.2 & 78.1 & \textbf{41.5} \\
\bottomrule
\end{tabular}
}
\caption{\textbf{Evaluation results.} For attribution, we report the AIS scores of the texts both before and after editing (before $\rightarrow$ after). For preservation, we report intent preservation $\presintent$, Levenshtein similarity $\preslev$, and the combined $\prescomb$. We summarize $\attrais$ and $\prescomb$ using their harmonic mean ($\HM{}$). %
}
\label{tab:main_results}
\end{table}


\begin{figure}[!t]
\centering
\includegraphics[width=\linewidth]{figures/fate_plot6tight.pdf}
\caption{
\label{fig:plot2d}
\textbf{Attribution and preservation scores.}
Dashed lines indicate the highest attribution score obtained by any of the models \emph{before} editing: points above the line have better attribution after revision.
The contours are $\HM{}$ level curves: points along a contour have equivalent $\HM{}$. Different models make very different trade-offs between attribution and preservation. Only \approach{} has a robust $\HM{}$ across all tasks.
}
\end{figure}


\subsection{Results}\label{sec:main-results}

For the main experiments, we report results on passages generated by PaLM and LaMDA. Results on GPT-3 passages show similar trends (Appendix~\ref{app:more-experiments}).
Table~\ref{tab:main_results} and Figure~\ref{fig:plot2d} show attribution and preservation results for each model and dataset. We also report $\HM{}$, the harmonic mean of the two metrics,
which is shown as level curves in Figure~\ref{fig:plot2d}.

\approach{} significantly improves attribution while preserving most of the original text.
In terms of $\HM{}$, \approach{} is the only method that performs robustly across all three datasets, and significantly outperforms prior methods on NQ and SQA.

We found that \approach{} is the only method that preserves the original intent of $x$ over 90\% of the time --- EFEC and LaMDA only manage to preserve the original intent 6--40\% of the time.
We also see that editing is crucial to improve attribution: if we only retrieve evidence to support the original response $x$ without editing, attribution ranges from the low 10s to mid 30s. After editing, \approach{} can increase attribution by up to 13\% absolute, while changing only 10--20\% of the text.

As noted in Section~\ref{sec:task},
one can  sacrifice preservation for higher attribution. EFEC is able to obtain strong $\HM{}$  on QReCC by making larger changes to the text in exchange for a higher attribution score. However, it occupies a very different point from \approach{} on the attribution-preservation trade-off curve, as visualized in Figure~\ref{fig:plot2d}.

\section{Analysis}\label{sec:analysis}

\subsection{Qualitative analysis}

\paragraph{Human oracle} To understand the remaining headroom in our task, we ask: \emph{what is the minimal amount of editing needed to make a text passage fully attributed?}
The answer would depend on the quality of the LM that generated the text as well as the task difficulty.
As an approximation, we manually edited 30 examples in our NQ benchmark until we judged them to be 100\% attributable. We achieved a preservation score of 88\%, which (when combined with 100\% attribution) translates to 93.6 $\HM{}$, indicating a significant headroom.

\paragraph{Analyzing the baselines}
\begin{figure}[!t]
\small
\centering
\adjustbox{max width=\columnwidth}{%
\begin{tabular}{
@{\!\!}r@{\hspace{3pt}}m{0.5\columnwidth}l@{}
}
\toprule

$x$: & \cmpR{Justice Ashok Kumar Mathur headed the 7th central pay
   commission in India. It was created in 2014 and submitted its
   report in {\color{red}\textbf{2016}}.} \\
& \textbf{Attribution:} 50\% & \textbf{Preservation:} 100\% \\
\midrule
EFEC: & \cmpR{The 7th central pay commission in India was created in 2014.} \\
& \textbf{Attribution:} 100\% & \textbf{Preservation:} 0\% \\
\midrule
LaMDA: & \cmpR{
I heard the 7th CPC made recommendations for
  increasing the minimum salary pay from Rs 7066 to 18k per month for
  new central government employees.} \\
& \textbf{Attribution:} 0\% & \textbf{Preservation:} 0\% \\
\midrule
\approach{}: & \cmpR{
Justice Ashok Kumar Mathur headed the 7th
central pay commission in India. It was created in 2014 and
submitted its report in {\color{blue}\textbf{2015}}.
} \\
& \textbf{Attribution:} 100\% & \textbf{Preservation:} 100\% \\
\midrule
evidence: & \cmpR{
The 7th Central Pay Commission (Chair: Justice A. K. Mathur)
\textbf{submitted its report on November 19, 2015}. The Commission had been
\textbf{appointed in February 2014}, to look at remuneration for central
government employees. \dots %
}
\\
\bottomrule
\end{tabular}
}
\caption{
\label{fig:compare_outputs}
\textbf{Example model outputs and human judgment of their attribution and preservation scores.}
EFEC reduces the passage $x$ into a single sentence. LaMDA changes the writing style.
\approach{} preserves the structure of the input passage. We show one evidence retrieved by \approach{} to help explain the example.
}
\end{figure}
As exemplified in Figure~\ref{fig:compare_outputs},
EFEC frequently attempts to summarize the entire passage into one sentence, or drops later sentences. This is likely due to EFEC's training data, which was limited to single sentences. This behavior generally increases the attribution score, because it is usually easier to make one sentence fully attributable than many sentences.
However, in datasets where the claim contains multiple sentences (NQ and SQA), such a behavior yields low preservation scores, and also results in outputs that are less informative.
We expect that EFEC could perform much better if its training data were augmented to include multiple sentences. LaMDA Research achieves similar attribution scores to \approach{}.
But as mentioned in Section~\ref{sec:baselines}, the intent and linguistic style of the output tend to deviate from the input, resulting in lower preservation scores (Figure~\ref{fig:compare_outputs}).
We emphasize that this is not a purely apples-to-apples comparison since LaMDA was not optimized for preservation.
Overall, these experiments are mainly meant to illustrate that prior models were simply not designed for the task of \taskName, rather than to mark \approach{} as the best method.


\paragraph{Analyzing \approach{}}

For the research stage,
the question generation model had comprehensive coverage:
a manual inspection of 40 examples shows > 80\% with questions that fully cover all aspects of the input text.
The retriever was strongest at researching content involving distinct entities (e.g., a movie, a major event, or a person). In contrast, we found significant headroom for better attribution of statements involving generic objects and more abstract claims (e.g. \emph{``Video games require electricity.''}--- since this is obvious to most humans, retrieved articles from the web tend to address related but different topics).
We suspect that a significant amount of attribution headroom on our benchmarks would benefit from a better  research stage.

\begin{figure}[!t]
\centering\small %
\adjustbox{max width=\columnwidth}{%
\begin{tabular}{@{}r@{\hspace{4pt}}p{1.25\columnwidth}@{}}
\toprule
\multicolumn{2}{c}{(a) \textbf{Correctly revising an entity}} \\
$y$: & If She Knew What She Wants was written by {\color{red}\textbf{Henry Roth}}. \\
$e$: & [en.wikipedia.org] ``If She Knew What She Wants'' is a song written by American singer-songwriter \textbf{Jules Shear} and introduced on \dots \\ %
$y'$: & If She Knew What She Wants was written by {\color{blue}\textbf{Jules Shear}}. \\
\midrule
\multicolumn{2}{c}{(b) \textbf{Correctly revising a number}} \\
$y$: & God Save the Queen became the British national anthem in {\color{red}\textbf{1745}}. \dots \\
$e$: & [www.britannica.com] The oldest national anthem is Great Britain's ``God Save the Queen,'' which was described as a national anthem in \textbf{1825}, \dots \\ %
$y'$: & God Save the Queen became the British national anthem in {\color{blue}\textbf{1825}}. \dots \\
& \emph{(The year 1745 was when the song was first performed.)} \\
\midrule
\multicolumn{2}{c}{(c) \textbf{Performing a necessary larger revision}} \\
$y$: & ``It's My Party'' is a song written and composed by American singer-songwriter and producer {\color{red}\textbf{Walter Gold}}. The song was published in {\color{red}\textbf{1962}} and first recorded by American singer-songwriter, {\color{red}\textbf{producer and arranger}} Lesley Gore. \\
$e$: & [secondhandsongs.com] Song: It's My Party written by \textbf{John Gluck, Wally Gold, Seymour Gottlieb, Herbert Weiner}\dots \\
$y'$: & ``It's My Party'' is a song written and composed by {\color{blue}\textbf{John Gluck, Wally Gold, Seymour Gottlieb and Herb Weiner}}. The song was published in {\color{blue}\textbf{1963}} and first recorded by American {\color{blue}\textbf{singer}} Lesley Gore. \\

\midrule
\multicolumn{2}{c}{(d) \textbf{Confused by misleading evidence}} \\
$y$: & In the Heat of the Night is a police procedural crime drama TV series produced by MGM Television. It aired on {\color{blue}\textbf{NBC}} from September 8, 1988 to May 3, 1992. \dots \\
$e$: & [metv.com] Watch In the Heat of the Night on \textbf{MeTV}. \\
$y'$: & In the Heat of the Night is a police procedural crime drama TV series produced by MGM Television. It aired on {\color{red}\textbf{MeTV}} from September 8, 1988 to May 3, 1992. \dots \\
& \emph{(The show did actually air on NBC in the 1980s, but the model fails to realize that MeTV is a re-run much later in time.)} \\
\midrule
\multicolumn{2}{c}{(e) \textbf{Improved attribution, but reasoning doesn't adapt to new info}} \\
& \textit{(Q: Does Homer Simpson need two hands worth of fingers to count to 5?)} \\
$y$: & Homer Simpson has {\color{red}\textbf{5}} fingers and 2 hands. Therefore, he does not need two hands worth of fingers to count to 5. He only need one hand to count to 5. \\
$e$: & [refinery29.com] Other iconic animated characters, such as Homer Simpson and his brood, \dots only have \textbf{four fingers}, too.\\
$y'$: & Homer Simpson has {\color{blue}\textbf{4}} fingers and 2 hands. Therefore, he does not need two hands worth of fingers to count to 5. He only need one hand to count to 5. \\
& \emph{(Having only  4 fingers per hand, he does need two hands to count to 5.)} \\

\bottomrule
\end{tabular}
}
\begin{subfigure}{0pt}\phantomsubcaption\label{fig:interesting_examples_entity}\end{subfigure}
\begin{subfigure}{0pt}\phantomsubcaption\label{fig:interesting_examples_number}\end{subfigure}
\begin{subfigure}{0pt}\phantomsubcaption\label{fig:interesting_examples_logic}\end{subfigure}
\begin{subfigure}{0pt}\phantomsubcaption\label{fig:interesting_examples_misleading}\end{subfigure}
\begin{subfigure}{0pt}\phantomsubcaption\label{fig:interesting_examples_reasoning_unchanged}\end{subfigure}
\caption{
\label{fig:interesting_examples}
\textbf{Example revisions  from \approach{}, both good and bad.}
$y$ = partially edited passage; $e$ = evidence; $y'$ = passage after editing with $e$.
}
\end{figure}

For the revision stage,
\approach{} was able to revise many unattributed claims, especially those involving entities and numbers (Figures~\ref{fig:interesting_examples_entity}~and~\ref{fig:interesting_examples_number}).
It can also perform larger revisions when necessary (Figure~\ref{fig:interesting_examples_logic}).
Moreover, \approach{} abstains from editing when the claim is already well-attributed:
on NQ, among the inputs with near-perfect attribution (pre-edit $\attrais{} > 0.9$), RARR does not make an edit in 90\% of the cases.
However, the system also has several shortcomings.
Some erroneous edits arise from %
misleading irrelevant evidence (Figure~\ref{fig:interesting_examples_misleading}).
We also observed an interesting challenge when revising reasoning chains, %
where the model successfully revised an incorrect claim, but did not revise subsequent reasoning steps that depend on the earlier claim (Figure~\ref{fig:interesting_examples_reasoning_unchanged}). In this case, further editing to improve logical coherence could help.

\subsection{Ablations} \label{sec:ablations}

\begin{table*}[!t]
\centering\small\setlength{\tabcolsep}{5pt}
\adjustbox{max width=\textwidth}{%
\begin{tabular}{
@{}l
r@{ $\to$ }rcc
r@{ $\to$ }rcc
r@{ $\to$ }rcc
@{}}
\toprule
 & \multicolumn{4}{c}{PaLM outputs on NQ}
 & \multicolumn{4}{c}{PaLM outputs on SQA}
 & \multicolumn{4}{c}{LaMDA outputs on QReCC}
 \\
 \cmidrule(lr){2-5} \cmidrule(lr){6-9} \cmidrule(lr){10-13}
 Model
 & \multicolumn{2}{c}{$\attrauto{}$} & $\preslev{}$ & \multicolumn{1}{c}{$\HM{}$}
 & \multicolumn{2}{c}{$\attrauto{}$} & $\preslev{}$ & \multicolumn{1}{c}{$\HM{}$}
 & \multicolumn{2}{c}{$\attrauto{}$} & $\preslev{}$ & \multicolumn{1}{c@{}}{$\HM{}$}
 \\
\midrule
Full \approach{}
& 45.6 & 54.9 & 89.6 & \textbf{68.1}
& 37.6 & 45.1 & 89.9 & 60.0
& 18.8 & 29.4 & 80.2 & \textbf{43.1}
\\ %
no agreement model
& 45.6 & 50.6 & 82.6 & 62.8
& 37.8 & 46.9 & 83.4 & 60.0
& 18.8 & 28.8 & 72.0 & 41.2
\\ %
query = input
& 45.4 & 47.2 & 98.4 & 63.8
& 39.4 & 30.3 & 98.8 & 46.4
& 19.7 & 20.6 & 96.3 & 34.0
\\
query = sentence
& 49.1 & 52.1 & 97.0 & 67.8
& 43.7 & 44.3 & 98.8 & \textbf{61.2}
& 19.0 & 19.6 & 97.0 & 32.6
\\
\bottomrule
\end{tabular}
}
\caption{\textbf{Ablation results.} We report the automatic metrics: $\attrauto$, $\preslev$, and harmonic mean between the two ($\HM{}$). We show auto-AIS scores both before and after editing (before $\rightarrow$ edit), with respect to the attribution report $A$ produced by the model.
Even though sentence-as-queries may achieve similar $\HM{}$ as \approach{}, it is less robust to corpus shifts and tends to retrieve passages that may encourage confirmation bias.
}
\label{tab:ablation_results_lower}
\end{table*}


\paragraph{Ablating query generation}
\approach{} uses generated questions as search queries for evidence retrieval. We consider two natural alternatives: using the entire input passage as a single search query, or using each sentence as a search query.
For the former, we retrieve $J = 3$ evidence snippets to make the amount a closer match to other methods.

The results are in Table~\ref{tab:ablation_results_lower}.
Using the entire input passage as the query gives poor results,
as the retrieved evidence tends to not focus on potentially unattributed parts in the passage.
Using sentences as queries gives results closer to the full CQGen, but a closer analysis reveals two caveats.

\begin{table}[!t]
\centering\small
\adjustbox{max width=\columnwidth}{%
\begin{tabular}{@{}lcccc@{}}
\toprule
 & \multicolumn{2}{c}{NQ $\HM{}$}
 & \multicolumn{2}{c}{SQA $\HM{}$}
 \\ \cmidrule(lr){2-3} \cmidrule(lr){4-5}
 Model
 & \multicolumn{1}{c}{orig} &  \multicolumn{1}{c}{no wiki}
 & \multicolumn{1}{c}{orig} &  \multicolumn{1}{c@{}}{no wiki} \\
\midrule
Full RARR
& \textbf{68.1} & \textbf{64.3}  & 60.0 & \textbf{57.6} 
\\
query = sentence
& 67.8 & 60.3 & \textbf{61.2} & 56.7
\\
\bottomrule
\end{tabular}
}
\caption{
\textbf{The impact of excluding Wikipedia from the retrieval corpus.} %
CQGen (full \approach{}) is more robust to Wikipedia's absence, while using sentences-as-queries suffers a bigger drop in performance.}
\label{tab:wiki_ablation_results}
\end{table}

First, sentences-as-queries are more effective when such sentences ``mimic'' content on the Web,
and are less effective otherwise.
In Table~\ref{tab:wiki_ablation_results}, we test this by excluding all of Wikipedia from web search results (since many PaLM outputs for NQ have a Wikipedia style). The attribution performance of sentences-as-queries drops significantly, while CQGen is more robust.

Second, sentence-as-queries tends to retrieve passages that may encourage confirmation bias. 
Consider the example \textit{``Georgia is called the Peach State, but California actually produces the most peaches.''} Retrieval using sentences-as-queries found an article echoing that California produces the most peaches, while CQGen generated the more impartial query \textit{``Which state produces the most peaches?''} and found a newer article saying that South Carolina replaced California as the top peach producer. In this case, \approach{} using CQGen needs to sacrifice more preservation score to edit the text, leading to a lower $\HM{}$ score. This underscores that attribution alone cannot measure ``correctness'' since not all evidence is up-to-date or reliable.

\paragraph{Ablating agreement model}
We try removing the agreement model, which effectively forces the model to revise the passage based on every retrieved evidence. The results are shown in Table~\ref{tab:ablation_results_lower}. As expected, more revision leads to less preservation score and spurious changes to the text passage, as demonstrated in Figure~\ref{fig:ablate_agreement}. %

\begin{figure}[!t]
\centering\footnotesize
\adjustbox{max width=\columnwidth}{
\begin{tabular}{@{}r@{\hspace{4pt}}p{1.2\columnwidth}@{}}
\toprule
$x$: & The Crown-of-thorns starfish is native to the Great Barrier Reef\dots The starfish was introduced to the Great-Barrier-Reef by \textbf{ocean currents}. \\
$e$: & [invasivespeciesinfo.gov] \textbf{Ballast water} is one of the major pathways for the introduction of nonindigenous marine species\dots \\
$y$: & The Crown-of-thorns starfish is native to the Great Barrier Reef\dots The starfish was introduced to the Great-Barrier-Reef by {\color{red}\textbf{ballast water}}. \\
\bottomrule
\end{tabular}
}
\caption{
\label{fig:ablate_agreement}
\textbf{Disabling the agreement model leads to over-edits.} Here, the evidence $e$ does not explicitly disagree with $x$, but without an agreement model to detect this, the edit model makes an unsupported change.
}
\end{figure}









\begin{figure}[!t]
\centering
\hspace*{-0.2cm}
\includegraphics[width=0.98\columnwidth]{figures/fate_e2e_barchart.png}
\caption{
\label{fig:e2e}
\textbf{Downstream task performance} on NQ and SQA. \approach{}'s revisions lead to better answer accuracy on NQ. No models improved answer accuracy on SQA. %
}
\end{figure}


\paragraph{Impact on downstream task performance}\label{sec:e2e_eval}
We have measured preservation using the metric defined in Section~\ref{sec:preservation}. However, another measure of preservation is whether the revised text can still be used to perform the task that it was originally generated for. 
Following EFEC,
we quantitatively evaluate this on short answer tasks NQ and SQA, and we summarize the result in Figure~\ref{fig:e2e}.

For NQ, each original text $x$ is a long-form response to a factoid question. To determine whether the revised text $y$ still serves this purpose, we feed the factoid question and $y$ back into PaLM and prompt it to extract a \emph{short} answer from $y$.
We find that \approach{} not only preserves the short answer accuracy but actually improves it by roughly 5\%.

For SQA, each original text is a reasoning chain that helps to answer a yes/no question. 
We feed the SQA question and $y$ back into PaLM and prompt it to output a yes/no answer, and evaluate answer accuracy. Here, we find that increasing attribution comes at a slight cost in downstream task performance: answer accuracy drops modestly for \emph{all} revision models (up to 2.6\%). 
We suspect that this may be due to noisy retrievals, which sometimes provide misleading evidence (exemplified in Figure~\ref{fig:interesting_examples_misleading}).
Furthermore, even though revisions can address factoid errors in the passage (e.g., \emph{``Homer Simpson has 5 fingers''} from Figure~\ref{fig:interesting_examples_reasoning_unchanged}), \approach{} currently does not try to modify subsequent reasoning steps which may no longer be logically entailed (e.g., \emph{``He only needs one hand to count to 5''}).











